\section{结论}\label{sec:conclusions}

我们给出了梯度流定义的能动量张量的普适 Wilson 系数到 \nnlo\ QCD。\nnlo\
修正使三个数值上占主导的系数 $c_1$、$c_2$、$\ringc_3$ 在中心标度
$\mu_0=3$\,GeV（$\mu_0=130$\,GeV）处改变约 10\%（1--2\%），其中 $\mu_0$
通过 \eqn{eq:mu0} 与流时间 $t$ 相联系。相对于把重正化标度绕其中心值变化
因子 2 所导出的 \nlo\ 结果，我们观察到理论不确定性的减小。第四个系数
$\ringc_4$ 的行为不那么令人满意，但其影响预期在数值上被压低。

除这一主要成果外，本文提出的新结果还包括：\msbar\ 方案下到 \nnlo\ 的流动
夸克场重正常数，以及构成 \emt\ 的常规 QCD 算符的反常量纲矩阵。

总之，我们希望这些结果有助于改进格点上 \emt\ 的研究。它们是梯度流形式内
高阶微扰计算系统化框架\,\cite{Artz:2019bpr}的第一个成果，该框架在这一理论
框架的其他应用中也应证明是有用的。
